\clearpage
\section*{Дополнительные материалы}
\addcontentsline{toc}{section}{Дополнительные материалы}
\setcounter{subsection}{0}
\renewcommand{\thesubsection}{S\arabic{subsection}}
\renewcommand{\thetable}{S\arabic{table}}
\setcounter{table}{0}

\subsection{Данные, определения показателей и чего не хватает}
\label{app:reproducibility}
Артефакт содержит записи отдельных выполнений для состояния модели и для
сравнения автоматических записей, окончательные агрегированные сводки
гибридного сравнения и итоговые сводки пилота окон. Все таблицы и
числовые утверждения формируются из этих источников. Гибридные сводки
содержат доли решённых и несовпадающие исходы, достаточные для разностей
и точных проверок Мак-Немара, но не траектории, промпты отдельных вызовов
или cache hit. Исходные сессии пилота недоступны; исходы по задачам
перенесены из журнала исследования.

Условные отношения ресурсов считаются на общих успехах:
\[
 \operatorname{median}_{(i,r):\,Y_{iar}=Y_{i0r}=1}
 \frac{X_{iar}}{X_{i0r}}.
\]
Нулевой знаменатель оставляет отношение неопределённым. Эффективность
использует все включённые исходы:
\[
 \text{Solved/\$}=\frac{\sum Y}{\sum D},
 \qquad
 \text{Solved/Mtok}=\frac{10^{6}\sum Y}{\sum(U+K+O)}.
\]
Solved/Mtok считает логические токены, а не равную работу для
кешированного и некешированного входа. Гибридные показатели эффективности
приближённые, поскольку суммы затрат округлены.

\subsection{Сопоставленные проверки исходов}
\label{app:contrasts}
\begin{table}[htbp]
\centering\small
\caption{Разности долей решённых задач в процентных пунктах с точными
проверками Мак-Немара. Бутстрэп-интервалы доступны только для состояния
модели.}
\label{tab:contrasts}
\input{../paper/tables/new-contrasts}
\end{table}

\begin{table}[htbp]
\centering\small
\caption{Разности долей решённых задач относительно полной истории внутри
когорт горизонта, заданных контролем (состояние модели).}
\label{tab:horizon-contrasts}
\input{../paper/tables/new-horizon-contrasts}
\end{table}

\subsection{Стоимость по всем сохранённым выполнениям}
\label{app:efficiency}
\begin{table}[htbp]
\centering\small
\caption{Стоимость и эффективность с учётом неудач. Медиана и среднее
стоимости отвечают на иные вопросы, чем суммарная стоимость на успех;
ни одна величина не оценивает стоимость повторов до успеха.}
\label{tab:economics}
\resizebox{\linewidth}{!}{\input{../paper/tables/new-economics}}
\end{table}

\subsection{Чувствительность стоимости автоматических записей к подмножеству успехов}
\label{app:composition}
Отношения на переменном подмножестве используют общие успехи каждой
конфигурации с полной историей; на фиксированном --- задачи, решённые
всеми конфигурациями. Богатейшие варианты стоят
\ExploreRichVariableCostRatio{} и \ExploreFailuresVariableCostRatio{}
от полной истории в первом случае, \ExploreRichFixedCostRatio{}
и \ExploreFailuresFixedCostRatio{} во втором: экономия относительно
ссылок на файлы сохраняется, относительно полной истории --- нет.
\begin{table}[htbp]
\centering\small
\caption{Чувствительность стоимости автоматических записей. Значения ---
медианы отношений к полной истории внутри задачи.}
\label{tab:composition-sensitivity}
\input{../paper/tables/new-composition-sensitivity}
\end{table}

\subsection{Отношения пилота окон по повторам}
\label{app:window}
\begin{table}[htbp]
\centering\small
\caption{Пилот окон: медианные отношения стоимости и пикового промпта
к полной истории в том же повторе, отдельно для каждого повтора.}
\label{tab:window-repeats}
\resizebox{\linewidth}{!}{\input{../paper/tables/new-window-repeat-sensitivity}}
\end{table}

\subsection{Сценарий кеш-скидки для исходных результатов}
\label{app:pricing}
Исходная публикация приводит только суммы токенов. Умножение её отношения
токенов полная история/состояние на предполагаемую кеш-скидку (hit полной
истории 0.93, состояния 0.72, цена кеша 0.198 от входа) даёт приведённый
ниже показатель. Он считает все токены кешируемым входом и служит
иллюстрацией чувствительности, а не реконструкцией счёта.
\begin{table}[htbp]
\centering\small
\caption{Опубликованное отношение токенов полная история/состояние
и показатель с учётом скидки по горизонту событий. Значения больше единицы
благоприятны для состояния.}
\label{tab:pricing}
\input{../paper/tables/new-pricing-sensitivity}
\end{table}
